% ===== §9 Eudaimonics: the field and the unfolding of one's own dynamics =====
\section{Toward a Eudaimonics of Relation: The Field and the Unfolding of One's Own Dynamics}
\label{p09:sec:eudaimonia}

\noindent
The keystone thesis tells us \emph{what} the good circle is---the non-possessive interpretive cycle that generates real value. It does not yet tell us how such a cycle is brought into being and lived. This is the threshold from the theoretical to the practical, and it opens onto a cluster of questions the paper has not so far confronted: How is the poetic language created and practised, such that the cycle becomes possible at all? How is the value of a system---of the two, of the family---to be perceived? How is value co-created? And how is the cycle made good, so that value returns to its creators? These are practical questions, and each merits its own sustained treatment; the present section does not pretend to answer them.\footnote{Each of these questions is, properly, the subject of its own study, and the series intends to take them up in turn after the present paper. What follows is not their resolution but the staking-out of the ground on which they will be worked, together with one organizing intuition---the field---that the author believes to be the right frame for all of them.} What it does is lay the ground for them, and offer one organizing intuition that the author believes to be the right frame in which all of them should be posed. That intuition is a \emb{field}; and the doctrine it opens is, in the end, a \emph{eudaimonics}---a theory of flourishing.

\subsection{From the presentation of a structure to the cultivation of a field}
\label{p09:sec:field-intuition}

Begin from a distinction that the whole of the foregoing has been pressing toward without naming. There are two radically different ways to conceive what it is one ``builds'' when one builds a good relation, or a good family, or a good community.

The first is to \emb{present a structure}: to lay down, in advance, a form---a set of roles, rules, scripts, an arrangement of who does what and stands where---and then to ask the participants to fill it. This is the natural mode of the engineer, of the planner, of the institution; and it is, as the foregoing sections have shown at length, the mode most prone to every pathology the paper has diagnosed. A presented structure is a settled structure: it declares, in advance, the shape the loop must take, and thereby forecloses the branching derivations that make a poem a poem. A presented structure is a script; and a script, however benevolent, assigns roles, and the assignment of roles is the first step toward reducing some participant to the executor of a part---toward, at the limit, the pure fuel.

The second is to \emb{cultivate a field}: not to lay down the form the participants must take, but to bring into being a \emph{condition} under which each participant's own, particular dynamics may unfold. This is the mode toward which the entire paper has been tending---it is non-possession given a positive, constructive name. Where the presented structure says ``here is the shape; take it,'' the cultivated field says ``here is a space whose curvature is such that what is in you may come out, and find, in coming out, that it is met.'' The field does not contain the trajectories; it shapes the space within which trajectories of their own accord arise.

\begin{claim}[The practical primitive is the field, not the structure]
The good relation is not a structure presented and filled but a field cultivated and inhabited. A structure prescribes trajectories; a field shapes the space within which each participant's own dynamics generate their own trajectories. The whole of the paper's negative wisdom---do not settle, do not siphon, do not fill the void with an idol, do not assign the script---is, positively stated, a single counsel: cultivate the field, and do not prescribe the path.
\end{claim}

The intuition has a precise resonance with the geometric language of the earlier sections, and the resonance is not idle. A geometric phase is the effect of a \emph{curvature}: the holonomy accrued around a loop is the integral of the curvature over the surface the loop bounds. The curvature is a property of the \emph{field}---of the connection on the bundle---and not of any particular path. To ``cultivate a field of positive curvature'' is, then, the exact geometric content of ``create the condition under which any path walked within it accrues a positive sublimation.'' One does not---cannot---hand a participant their sublimation; one can only shape the field so that, whatever path that participant walks of their own dynamics, the walking sublimates. \emb{This is the geometric meaning of \emph{wu wei}: not to draw the trajectory, but to tend the curvature; and then to let each go their own way, trusting the field.}

\subsection{Why the field and not the structure: the uniqueness of each participant's dynamics}
\label{p09:sec:why-field}

Why must the practical primitive be the field? The answer lies in a fact the structure cannot accommodate and the field is built to honour: \emb{each participant has their own, irreducibly particular dynamics.} No two persons sublimate the same real along the same path; no two encircle the same void in the same orbit. The structure, because it prescribes the path in advance, must presuppose a generic participant---an interchangeable filler of a role---and to the precise extent that a real participant departs from that generic template (which is to say, to the precise extent that they are a particular person at all), the structure does violence to them, forcing their singular dynamics into a path not theirs. The field presupposes no generic participant; it is, on the contrary, the condition under which the \emph{particularity} of each can come forth.

This is why the doctrine that the field opens is, at bottom, a \emph{eudaimonics}. For what is flourishing? Not the attainment of a prescribed end-state, not the successful filling of a good role, but---in a formulation the field makes available---\emb{the unimpeded unfolding, within a field that meets it, of one's own particular dynamics}. The Aristotelian \emph{eudaimonia} \parencite{aristotle2000} named the activity of the soul in accordance with its own virtue, the actualization of a being's characteristic potential; the Spinozan \emph{conatus} \parencite{spinoza1994} named each thing's striving to persevere in and express its own essence; the Daoist \zh{自然}---``self-so,'' the spontaneous---named the way of a thing when nothing forces it. Each of these, in its own register, names the same intuition the field formalizes: that flourishing is not conformity to an imposed form but the coming-forth of what a being, of itself, has it in it to become.

\begin{claim}[Flourishing as the unfolding of one's own dynamics in a meeting field]
Flourishing (\emph{eudaimonia}) is the unimpeded unfolding of a participant's own, particular dynamics within a field that meets them---a field of such curvature that the path they walk of themselves accrues a positive sublimation, and that sublimation refluxes to them. It is neither the attainment of a prescribed end (that is the settled structure) nor a private self-realization indifferent to others (that is the monad, which this series has from the first denied). It is relational and generative: one flourishes \emph{in} a field one also helps to cultivate, by unfolding dynamics that, in unfolding, tend the field for others.
\end{claim}

Two features of this formulation must be marked, for they distinguish it from the doctrines it draws on. First, it is \emph{relational}: the unfolding is not a solitary self-actualization but a coming-forth ``within a field that meets it''---and a field is never one's own alone but the joint work of all who inhabit it. The monadic eudaimonia, the flourishing of the self-sufficient individual, is foreclosed by the relational ontology of the series: there is no self-prior-to-relation whose potential could unfold in a vacuum. One's own dynamics are themselves generated in the field, even as they help generate it. Second, it is \emph{generative} in the precise sense the paper has built: the unfolding of one participant's dynamics, done within a good field, is itself a tending of the field's curvature for the others---the very act of one's flourishing cultivates the condition of another's. This is co-authorship (\xref{p09:sec:coupled}) seen from the side of the participant rather than the system: to co-author the poem is, for each, to flourish by unfolding one's own dynamics in a way that, in unfolding, makes the field more able to meet the rest.

\subsection{Happiness and value: a first principle}
\label{p09:sec:happiness-value}

This lets us state, if only as a first principle to be developed, the relation between flourishing and value that the foregoing theory of value (\xref{p09:sec:value}) implies. The connection runs through the three strata.

A flourishing built upon \emph{imaginary value} is the flourishing of the idol-encircler: intense, and unsustainable, for it depends upon an idealization that the real other must, in time, pierce. It is the happiness of the rose at the height of its bloom---resplendent because exhaustive. A flourishing built upon \emph{symbolic value} is the flourishing of accumulation and recognition: real, not to be despised, but exposed to siphoning and to inflation, and never, by itself, sufficient---for the symbolic can always be stolen, devalued, or generated cheaply, and a happiness resting on it alone rests on sand. A flourishing built upon \emph{real value} is the only one that can bear the good's phase from within: it is the happiness of encircling, with another, an unfillable void---of being witnessed, and witnessing, in one's irreducible particularity---and it does not exhaust, for it possesses nothing, and what is not possessed cannot be used up.

\begin{claim}[The first principle of a eudaimonics of value]
Sustainable flourishing is the unfolding of one's own dynamics in a field whose circulating value is led by \emph{real value}---the unpossessable, unsiphonable value generated only in encircling witness. Flourishing on imaginary value is resplendent and exhaustive (the rose); flourishing on symbolic value is real but insufficient and stealable (the accumulation); flourishing on real value is the one that sustains, because it possesses nothing and so uses nothing up. The cultivation of a field for flourishing is therefore the cultivation of a value-cycle led by real value: a field in which each, by unfolding their own dynamics, encircles with the others a void none fills, and the sublimation refluxes to each.
\end{claim}

This principle gives the four practical questions with which the section opened their common form, and shows why they are one question seen from four sides. ``How is the poetic language created and practised?''---it is the cultivation of a field of positive curvature, in which signs anchor without settling and invite without prescribing (\xref{p09:sec:semiotics}). ``How is the value of a system perceived?''---it is the discernment, beneath the imaginary and the symbolic, of the real value being generated in encircling witness, which is the labour of diagnostic attention the epistemology demanded. ``How is value co-created?''---it is the co-authorship by which each participant, unfolding their own dynamics within the field, generates the real value that the field, in turn, refluxes to all. ``How is the cycle made good, so that value returns to its creators?''---it is the unceasing interrogation of the pure-fuel criterion (\xref{p09:sec:fuel}), the vigilance that no participant's flourishing is the fuel of another's. Four questions; one field; one eudaimonics.

\subsection{A note on what remains}
\label{p09:sec:remains}

It must be said plainly that this section has done no more than name a frame and state a first principle. The hard work---how, concretely, a field of positive curvature is cultivated; how the particular dynamics of a participant are to be perceived and met rather than prescribed to; how the perception of a system's real value is to be practised against the constant pull of the imaginary and the symbolic; how co-creation is organized without lapsing into the assignment of roles; how the return of value to its creators is secured under real conditions of power and pressure---all of this lies ahead, and each piece of it will require its own patient development.\footnote{The author marks these as the practical sequel to the present paper: a eudaimonics of the cultivated field, a phenomenology of the perception of relational value, and an ethics of co-creation under asymmetry. The present section plants them; it does not pretend to have worked them.} The paper has, in this, been faithful to its own counsel: it does not settle what has not been walked. To present here a finished doctrine of practice would be to commit, in the very section that warns against it, a settling failure---to cash out, with a sign, a sublimation the work has not yet earned. The field, like the poem, is cultivated and not declared; and a theory of the field can do no more, at this stage, than tend the ground in which the later, more patient work may grow.
